\documentclass[11pt]{article}
\usepackage{../shared/luxcover}
\usepackage{amsmath,amssymb}
\usepackage{hyperref}
\usepackage{booktabs}
\usepackage{enumitem}
\usepackage{xcolor}

\title{Lux Platform}
\author{Zach Kelling\\\textit{Lux Industries, Inc.}\\\texttt{research@lux.network}}
\date{April 2024}

\begin{document}

\luxcoverpage

\begin{abstract}
This paper presents the design and architecture of the described system.
\end{abstract}

\tableofcontents
\newpage


\section{Abstract}

The current FinTech landscape is fragmented between banking and blockchain protocols, exchanges, and varying compliance and regulatory regimes. Decentralized consensus is needed to safely navigate this perilous landscape. Tools that unite various banking or blockchain networks while allowing frictionless movement between them are necessary. In addition, these tools must incorporate heretofore unbuilt functionality that allow them to function in a fully legal and compliant fashion.


Decentralized networks mechanize the trust function and obviate the need for human controlled centers of trust, creating new and efficient means of large-scale organization and cooperation. One of the most obvious applications is to the financial sector--the original ``trust center''. There have been many promises in this regard, but little action due to the myriad of problems that need to be solved. Security, fees, compliance, disparate regulation and technology are just a few of the issues that must be considered in order for mass adoption of blockchain technology by the financial sector.


Lux Finance is designed to address the needs of the financial services sector while also solving critical usability, scalability, and interoperability problems that are standing in the way of wide-scale adoption.


The Lux Finance Platform will be a solution that combines the highest standards in underlying security structure while employing a revolutionary wrapper that is agnostic of protocol and exchange. It will be able to contain all types of securities, collateral and other assets. The Lux Platform will finally permit tokenized assets to be natively swapped between different blockchains at de minimis costs and in a fully decentralized, yet regulatorily-complaint fashion. The applications are nearly infinite, including stable tokens, investment vehicles, all manner of assets, and a universal connective mesh that unifies the blockchain ecosystem.


\section{Lux Finance Platform}

\subsection{Problem with Current Solutions}

The current blockchain environment is revolutionary in certain ways, and limited in others. On the positive side, there are the first pulses of a truly decentralized system where the value is secured and held by a network of holders. However, on the other side, there is a user environment accustomed to a certain level of regulation and oversight. This is interesting and treacherous ground to cover.  How do you bridge the gap from a user base and society that is used to having their value secured by a third party?


The first problem is that of exchanges. In the world of blockchain, the most critical point of failure is that of an exchange. There is the danger of your funds and identity being compromised. You risk of loss of funds. There is an opaque order book. The fees are not transparent and obvious, especially when considering the promise of blockchain being transparent, decentralized and a driver in cost reduction.


An additional layer of concern is around issues of compliance. In the real world, where safety and compliance are secured through trusted institutions and regulatory bodies answerable to governments and people, security and safety have been established over a long and iterative process. Blockchain specifically promises to eliminate these areas of cost and potential failure. But how to bridge the gap from the current reality to the uncertain decentralized future?


We believe the answer is to embrace the role that regulation has played and understand, that prior to blockchain, these rights have been secured by regulatory bodies. It must be understood that the same security that people have come to expect, is the same environment this new technology must exist under. It is this expectation of a regulated environment that has made the public susceptible to unscrupulous deployment as this structure as a way to gather funding.


So the problem is multifaceted; a public that is used to trust and security being assured by a centralized power and the tools that allow compliant, transparent regulated products to be offered.


The Lux Platform aspires to meet both needs: providing the tools to offer decentralized products that are compliant and work with the current regulatory bodies to satisfy the current and appropriate level of government oversight. While the ultimate goal of blockchain technology is to be truly decentralized networks, the reality is that it must, at least, live up to the safety and security of the centralized systems it is replacing.


\subsection{The Exchange Problem}

The current blockchain ecosystem relies on a network of unaffiliated exchanges that allow people to move between various blockchains and tokens, as well as in-and-out of various real world stores of value such as currencies. There are problems with these exchanges including high and opaque fee structure, murky order books, and a dubious criteria for listing on exchanges. There is also the problem of decentralization; the antithesis of blockchains existence.  Does it have to be this way?


We Say No emphatically! Many of the innovations and issues that arise from the development of blockchain solutions are informed and warped by real world solutions. Blockchain was designed to solve the problems associated with centralization, but we often turn to centralization as the way to solve these problems because it is familiar and it is easier to implement. The exchanges are an excellent example of this.


We are accustomed to going to center of exchange for the transformation of dissimilar objects. We are accustomed to paying a fee for this and outsourcing the mechanism of this transaction to a third party. Blockchain and its innate trust and lowest cost nature is made for this.  Yet, the centralized, expensive and opaque exchange structure exists today. It does not have to and the Lux platform executes on this function at its most basic and at the lowest cost.


\subsection{Compliance Challenges}

Compliance, regulatory oversight and control are often considered incompatible in blockchain and cryptocurrencies. Many consider their advent the harbinger and destruction of these centralized areas of control. But one must ask themselves why were these regulations and systems set in place initially.


These rules and regulations were designed to secure the most trust dependant of networks. The issue is the tools available to the current blockchain do not even allow for the implementation of these regulations. The Lux Platform is designed specifically with these in mind and to be upgradeable for future innovations.


\subsection{Cost}

The elimination of cost and transaction friction is one of the aspects that blockchains and smart contracts are supposed to achieve. The reality however, is that every point of the ecosystem is layered with costs that far exceed the current centralized regime. The Lux Platform is capable of bringing these costs down to their lowest structural point.


\subsection{Gateway Between the Physical and Digital}

Anybody that has interacted with cryptocurrency or digital ledgers understands how important and tenuous their relationship is. For the most part, all of our real world costs are settled in fiat currencies. Any friction, complication, delay our cost is perceived as fairly fatal in any consumer adoption and use of cryptocurrency more broadly. And yet interest and utilization grows.


The Lux Platform aims to be the most cost effective, platform agnostic, fastest way to move between the physical world, whether fiat currency or other traditional securities. Many of the choke points are around Know Your Customer/ Anti Money Laundering (KYC/AML) requirements. These functionalities, among other compliance necessities are built into the Lux Platform. Solving these problems are some of the limitations that prevent blockchain technology from being adopted at scale and the consistency of the traditional finance system. Lux addresses these challenges.


\subsection{Interoperability Between Protocols}

One of the major benefits to blockchain technology is its decentralization; but this decentralization presents problems of its own. We have already seen the proliferation of systemically important protocols and tokens, but there has not been a way for them to interact with each other in a frictionless and cost effective environment. The protocol's independence is of central importance, but if they cannot work together how useful is it to broader society?


Currently the interplay between blockchains is handled in a very centralized, status quo method.  They are treated as individual and distinct entities to be exchanged like currencies, securities and other primitive devices with their attendant middlemen and costs. But the promise of this technology is supposed to spread the the center of trust and drive down the attendant cost of securing it. Independant, isolated blockchains, no matter how innovative their feature set are held hostage by the current connective mechanism: the exchange. The Lux Platform facilitates this interchange at scale and at the lowest possible cost, determined by the underlying protocol. It allows these innovations to interact with each other, in the manner that programs interact, at a basic level, in the most efficient way possible at the lowest cost.


\subsection{Technological Innovations}

Fully decentralized, non-custodial securities platform


Assets are never in our possession


Blockchain agnostic protocols


Atomic cross-chain swap protocol


Support for multiple chains and liquidity pools


Compliance Innovations


Notifications as an example


Redemptions


Shared  KYC/AML Registry


Regulatory controls


Versioned and upgradeable protocols


Easily adapt to changing compliance / regulatory needs


Quickly adopt rapidly developing protocol standards


Scaling via ZK Proofs


Execute smart contracts off-chain, blockchain enforces smart contract code only in case of disagreement


Improved privacy, scalability


Oracles


Connect real world data and third-party systems with smart contract code


Automated interest disbursements


DEX


Securities focused DEX with real-time transactions


\subsection{Applications}

The applications of this technology are quite broad. To provide an interconnectivity between blockchains with the attendant compliance functionality has myriad applications. The Lux Platform focuses on the challenges and needs presented by the regulated financial services industry. These capabilities will be applicable to the broader world, but are immediately necessary and innovative to financial services.


\subsection{Stability}

The world is awash in stable tokens in various mechanisms to maintain stability, but often overlook the most fundamental functionality. To be truly useful, stable tokens most me liquid and as fungible as possible. Linking it to an underlying store of value is a secondary function. All current stable tokens are organised around the dubious current exchange ecosystem. None are exchangeable across multiple blockchains and are nearly impossible to redeem. Additionally, they have not been designed to be compliant and legal. This frictionless structure, coupled with the underlying compliantly represented security, offers the best, most liquid and flexible stable token solution. Additionally, the fee transparency and seamless nature of the Lux Securities Platform allows the benefits and drawbacks of the underlying protocol to present themselves.


\subsection{Universal Means of Exchange}

Many blockchains aspire to be the means of exchange of this new ecosystem. None has emerged as an obvious leader or seized a dominant, unassailable position. Why?


The functionality, cost, stability and interoperability has not been realized yet. The Lux Platform has the requisite feature set and potential ability to operate at that basic layer to fill that role.


\subsection{Near Zero Fee Products}

By lowering fees at the protocol layer and allowing liquid, compliant flexible ownership of the underlying asset, these products are finally able to deliver on one of blockchains' ultimate promises: the reduction of fees through automation and the elimination of the cost of centralizing trust. To date, the solutions have been orders of magnitudes more expensive than the traditional financial service products. The Lux Platform allows for the deployment of the structurally lowest fee structure of financial service products possible and finally realizing the ambition of efficient low cost blockchain products


\subsection{The Connective Tissue of the Blockchain Space}

Until now, the blockchain space has been a collection of competing and unrelated protocol. The vision of fast, cheap decentralized solutions to our current centralized world has not materialized. While the current financial system is centralized, it is also fast and reliable. This is accomplished by agreed upon communication and security protocols like SWIFT. How do we achieve such harmony and communication between the ever evolving and independant areas of the blockchain ecosystem? The Lux Platform offers a potential solution to this basic and critical functionality. By offering this connective piece at the lowest possible cost in both time and speed, while being protocol andgostic, Lux has provided the universal underpinning that is necessary for this concept to reach the diversity scale that is possible.


\section{Lux Treasury Token}

\subsubsection{Underlying Structure}

Closed End. Legally, the Lux stablecoin will be digitized shares of a registered investment company formed pursuant to the Investment Company Act of 1940.  However, unlike traditional mutual funds, the Lux stablecoin will be closed-end, with regular and frequent redemption periods.  We anticipate the legal structure will be sufficiently robust to support peer-to-peer transactions as well as transactions on exchanges.  Since the underlying assets will be actual US Treasury securities, which are very liquid, the fund will be able to meet any redemption requests to exchange the shares for fiat.  All such requests will be processed through a registered custodian (such as a U.S. bank) and wired directly to the bank account of the investor.


\subsubsection{Fully Compliant and Regulated}

The legal structure will be fully compliant with, and registered under, U.S. securities laws and subject to oversight by the SEC.  In fact, before the stablecoin can be offered to investors, the SEC must clear the registration.  Once the stablecoin fund is fully registered, it will be available for sale for all investors in the U.S., regardless of accreditation status.  Further, since U.S. securities laws are generally considered the most stringent in the world, a registered U.S. offering would generally be a permissible investment for all retail investors in other jurisdictions outside the U.S. with minimal, if any, additional regulatory pre-conditions.


As an Investment Company Act fund, the fund will be required to provide periodic statements of performance and holdings to investors.  In addition, the fund will be subject to various investor protections such as the requirement to hold custody of assets with a qualified custodian, the requirement to engage a qualified transfer agent to keep track of shareholders, and the requirement to perform semi-annual and annual audits by a PCAOB auditor (i.e., an auditor approved by the SEC).  This assures investors that their stablecoins are fully-backed by investments held in institutions qualified to ensure their safety, and offers a level of transparency currently lacking in other stablecoin projects.


\subsection{Potential Underlying Holdings for Lux Platform}

\subsubsection{Lux Dollar, US Treasuries backed stablecoin}

The Lux Platform will launch with the Lux Stablecoin, a financial service product backed by US Short-term US treasuries. While marketed and intended to fill the need of a truly universal, cheap liquid, and compliant stablecoin, its underlying security adds value as an investment vehicle not offered by current solutions. By holding of US treasuries, with their unmatched liquidity and regulatory preference, matched with their actual yield, transforms the nature of stabletokens. The current landscape of stabletokens is populated by unregulated, unredeemable, possibly illegal structures that are expensive and are only popular due to the lack of better alternatives.


Current projects that focus heavily on the stability mechanism, seem to lack rigor around the legal structure and how they will practically work in this space. The Lux stable token will be the most liquid, universal, free to consumer stable token in existence when deployed.


\subsubsection{Puerto Rican Bonds}

To address another specific objective of interest to the crypto community, the second offering on the Lux Platform road map is a Puerto Rican Bond Portfolio Token. The blockchain community has shown great interest in Puerto Rico, both as a  area that can benefit from blockchain community investment as well as a domicile that is tax advantaged for residents. This offering highlights the flexibility of the Lux Platform while employing marketable end products.


\subsubsection{Anything that is Currently Held In Traditional Securities Wrappers}

The first two Lux Platform products are intended to highlight the benefits of the structure while addressing specific needs of the current, albeit small, blockchain community. There is nothing limiting this vehicle's application to wrap other securities and it is Lux's intention, to launch, market and distribute these products to meet the needs of all inventors, no matter their return goals and risk profiles.


\subsection{Resources Needed to Succeed}

To execute on this broad and ambitious goal, expertise is needed in several specialized and unrelated areas. To date, attempts to solve these problems have often been strong in one or two areas, but not necessarily in all the areas. The difficulty of arranging world-class resources in all these areas and having them come together to make a integrated products of this scale and complexity is quite difficult. The difficult nature of this project mirrors the challenges faced by the broader blockchain space as the technology itself often requires high levels resources in these unrelated areas. This fragmented approach has led to some of the confusion and disunity in the space, especially in the area of financial services. For instance, technological innovation has not always contemplated the regulatory and legal ramifications of these products and has led to solutions that often run afoul of the law.


The Lux Platform has been assembled with a world class resources in all the necessary areas and tightly integrated them.


\subsection{Technological Expertise}

Technological expertise is obviously necessary, and is rarely overlooked, but the particular type that is necessary is often misunderstood. As blockchain technology is so new and evolving so quickly, it has been hard to determine what capabilities are vital to create these products. We view that experience in securities, cryptography consensus networks and building systems at scale are vital. Again, financial service blockchain products are securing people's money in a decentralized, public facing way the before this had been done in a highly gated, secure and closed networks. This presents a unique set of challenges that have not had long periods of deployment to test.


\subsection{Structuring and Legal Expertise}

Often overlooked on the blockchain product side is an expertise in structuring legal, compliant products that meet the high standard of the traditional financial service industry. Current solutions often make sacrifices in areas of functionality that would not be tolerated by issuers or consumers of traditional financial service products.


Also, failure to understand how current securities laws apply to this new technology have exposed both blockchain companies and their customers to legal jeopardy. Any products offered in this space have to cdesigned applying this most stringent standards of both functionality and legal issues are to be overcome. This is vital for the technology to achieve its goal of mass adoption.


\subsection{Financial Service Expertise}

As the Lux Platform is both a financial services as well as technological solution, financial services expertise is necessary. Even though there is a ultimate goal of a completely decentralized world where products are deployed effortlessly and immediately adopted by consumers, this is not traditionally the way financial service products occur. There is knowledge of brand, consumer needs distribution to name just a few. Thes skills must be in place to design market and distribute products if they are to achieve mass adoption, especially in this period where the vast majority of assets and resources reside on the traditional side of the fence.


\subsection{Operational Expertise}

The scope and breadth of blockchain's possibility to revolutionize aspects of our society has generated tremendous enthusiasm and interest in the space. Additionally, the technology has made it possible to raise capital for projects in a much faster and often less scrutinized manner then the traditional methods that have preceded it (VC Fundraising etc) The disruptive nature of this is both a blessing and a curse. Those traditional gatekeepers, while centralized played an important role in gauging teams abilities to execute based on their current set of skills as well as their track records for execution. Blockchains generally and the ICO process specifically has eliminated that oversight function and have given teams with no history of execution. It is no surprise that there is a wave of unfinished projects and failed ICOs that are sweeping through the ecosystem. These are complex projects that require expertise in execution. While a history of past executive abilities is no guarantee of future success, it is an important data point when evaluating the probability of a project of this level of complexity being brought successfully to fruition.


\subsection{Lux's Combination of These Resource}

The Lux team has the above mentioned resources and capabilities. They are being deployed in an integrated fashion that is necessary for a project as technologically and legally complex as the Lux platform to succeed.


\section{Lux Network}

Interplanetary Commerce


What does commerce look like in 100 years


What is necessary at planetary and interplanetary scale


What system is necessary to support billions of users


Own all the interconnect layers


Network token pays for all fees


DNS resolutions


Decentralized computation marketplace


Transactions


Services at service layer providing


Governance


Messaging


Notifications


Payments


Exchange


We want to become the Apache web server of the blockchain generation


Flexible building blocks


Stellar Consensus Protocol 2.0


Outdated tech, right ideas


Decentralized federation


Mainnet is just our net


Free form path finding


Each chain has own governance / hardware


Google-scale decentralized companies


Decentralized internet 3.0


Satellite network


Appropriate for moon and beyond


Latency compensation


Immortal


Byzantine-fault tolerant


Support for massive network loss and partitioning


Push notifications


Quantum computing proof


Fee structures


Network token for blockchain platform


Fees per transaction / generally subsumed by businesses


Free for token holders


Should be able to use any natively minted token as payment based on built-in exchange


\subsection{Architecture / Design}

Block DAG architecture


Designed for SSDs (LSM-tree KV store, Graph DB per node)


Distributed / Highly available


Sharded storage


DPoS / SCP federated consensus


Pluggable architecture


Customize computation engine, pluggable smart contract implementations


Initially own all hardware


Enable anyone to participate, license to larger institutions


Managed data pipes, guaranteed bandwidth / scalability


Service / maint by us


Forced upgrades


Guaranteed performance levels / hardware standards


Institutional grade compliance / security


New levels of compliance / security and performance standards / metrics


Highest performance / scalability


Millions of concurrent real-time users available


Multi-tier high performance caching layers


Flexible / Modular design / developer-friendly framework / toolkit


Multiple industry support


Open source components, DIY blockchain


Closed-source, highly scalable on-premise or self service solutions


\subsection{Consensus Protocols}

\subsubsection{Delegated Proof of Stake}

Emerging as industry standard / Easy to replicate EOS


Staking enables voting for BPs / plutocratic governance


Evolve to support decentralized identity / 1 vote per wallet per human


\subsubsection{Proof of Reputation}

Peer to peer decentralized reputation tracking


Kevin Baconized so we can back trace who trusts who and evaluate and rank trust structures


This is similar to what Stellar kind of does right now.


We need to figure out whether or not a trust endorsement matters more if someone is highly trusted themselves (aggregate trust)


Blocklists by consensus to remove bad actors and spammers


\subsubsection{Proof of Health}

Consensus algorithm automatically audits the integrity of the node


Blockchain analytics tools built in at lowest level (see below)


\subsubsection{Proof of Service/Replication/Availability CDN}

CDN layers will utilize proof of storage to pay people for storing content


Novel algorithm for evaluating how quickly content is served and what content is stored based on availability parameters


\subsubsection{Sidechains / Sharding / Federated Consensus / Path finding}

Sidechain per app


Should go to the point of having sidechain primitives in dev model


In the same general way golang as threading primitives


\subsection{Hosting / Network operation}

Distributed DNS for naming nodes in the blockchain


\subsection{Decentralized computation}

Cloud Function and Smart Contract Library/Package Manager


Allow easy distribution of libraries to accelerate development and adoption


Support Go, Node, Python, EVM


\subsection{Decentralized Oracle service}

Ability to write / read from blockchain with decentralized trust


Watcher nodes incentivized by network token


\subsection{Blockchain-as-a-Service / Consulting / Services}

Hosted blockchain nodes


Ability to scale / launch nodes for multiple chains


We develop your financial product as a service


Use our products for other companies, institutions


\section{Lux Security Token}

\subsubsection{Applications}

Stable tokens


STOs, shares


Real-world assets


Real estate


Fungible and non-fungible


Indexes


Funds


\subsubsection{Protocol}

Native ERC-20 / EOSIO token support


Works in any wallet


Restricts transfer to only between compliant addresses


Compliance / Regulatory features


Restricted trading based on KYC/AML, compliance, regulator


Decentralised registry network providing trust network / decentralized regulatory approval / KYC / AML


Token economy controls


Ability to freeze and migrate across chains


Ability to blacklist bad actors / freeze funds


Full version control / upgradeability


Data layout can never change (data will never be lost)


But new functionality can be layered on top


Compliance / Regulatory controls pluggable / upgradeable


Private chain support


Full off-chain persistence and legacy interoperability


Completely decentralized asset storage / accounting


\section{Lux Securities Platform}

A single platform for all securities issuance and trading


Registry and issuance of securities


Next-generation blockchain agnostic securities framework


Support for multiple securities protocols


Lux ST


Polymath ST-20


Harbor R-Token


TrustToken


etc...


Built-in wallet with full support for push notifications and securities tokens


Cross-chain atomic swaps


Built-in DEX


KYC / AML


\section{Lux Wallet}

Fiat on-ram/off-ramp


On-ramp:


Crypto payment / Fiat payment -> STBL / other crypto


Off-ramp:


Auto-generate debit card numbers for each transaction


More secure


Less prone to attack


Split payments


Bill pay with crypto


Scheduling / pooling for large payments


Support owned by Bot system (near human level of support)


\section{Lux Bridge}

Provides cross-blockchain liquidity


Supports security tokens and all compliance features


As long as users are KYC/AML'd with one of our registry providers with a trust line, we allow trades to happen


Registries automatically updated


\section{Lux Web3 Platform}

Ability to deploy cross-chain dapps and protocols at once


Ability to easily handle payments


Handle identity / authentication


Rapid scalable dapp development


\subsubsection{Developer Friendly}

Flexible object model


Automatic API generation


Multi-lingual SDKs


Bounty system


Drive open source adoption and community building


\subsubsection{Blockchain acceleration layer}

Internet existence services


DNS


Self-name node capability


Consensus used to describe what exists where


CDN


Redundant failover capability


Clustering capability


1M / TPS


Hosted NoSQL backend, fully query support


HTTP/HTTPS REST API


GraphQL \& SQL API Layer


Allow full drop in support for existing software drivers


Data Streaming API (WebSockets)


Push Notifications


Atomic cache-invalidation


\subsubsection{Blockchain analytics tools}

Publicly accessible health markers on network


Allow QOS


Low-level support for diagnostics on token ecologies and transactions


\subsubsection{Scalability via ZKP}

1M / TPS


Cross-chain Atomic Swap Layer


Ability to launch cross-chain tokens


EOS, ETH, XLM simultaneously


\subsubsection{Encryption Layer}

Configurable security contexts, encrypted yet shared data


Need ability to purge data


Off-blockchain data storage


Full near-realtime storage of chain state


Ability to use ZK proofs to interact with public chain data in a way that's fully anonymized


\subsubsection{Decentralized Message Queue / Event system}

Connective tissue for cross-chain operations


Pluggable second layer solution over individual chains


Allows developers and businesses to transcend a specific chain and remain structure agnostic


\subsubsection{Decentralized Cloud Functions}

Flexible programming model


Marketplace for decentralized computation


Arbitrary tags allowing hardware traits, usable inside programmable modules


\subsubsection{Secure / Trusted Computing Environment}

Trusted Platform Module (TPM)


Use latest advancements in security to execute smart contracts in Intel SGX / trusted computing hardware modules


Tiers of trust, highest layer will require TPMs to verify hardware


Completely isolated execution context


Complete key security


\subsubsection{Decentralized machine learning / AL Layer}

GPU-Accelerated nodes


High security specialized cryptography hardware (talk to Nvidia?)


\subsubsection{Quantum Computing Ready}

Maybe some quantum secured stuff with D-Wave systems


Just get a D-Wave system to try and crack the security?


\section{Research}

https://scalingbitcoin.org/milan2016/presentations/Scaling\%20Bitcoin\%203\%20-\%20Mark\%20Friedenbach.pdf


https://www.tik.ee.ethz.ch/file/a20a865ce40d40c8f942cf206a7cba96/Scalable\_Funding\_Of\_Blockchain\_Micropayment\_Networks\%20(1).pdf


\subsection{https://blockstream.com/bitcoin17-final41.pdf}

https://scalingbitcoin.org/milan2016/presentations/D1\%20-\%204\%20-\%20Andrew\%20Poelstra.pdf


\subsection{https://eprint.iacr.org/2016/1159.pdf}

\subsection{https://eprint.iacr.org/2018/104.pdf}

\subsection{https://github.com/ethereum/wiki/wiki/Sharding-FAQs}

\subsection{https://github.com/ethereum/wiki/wiki/Governance-compendium}

\subsection{https://tendermint.com/static/docs/tendermint.pdf}

\subsection{https://eprint.iacr.org/2017/634.pdf}

\subsection{https://cryptojedi.org/papers/dilithium-20170627.pdf}

\subsection{https://www.stellar.org/papers/stellar-consensus-protocol.pdf}

\subsection{https://bravenewcoin.com/assets/Whitepapers/QRL-whitepaper.pdf}

John R. Douceur, The Sybil Attack, International Workshop on Peer-to-Peer Systems (2002), pp. 251-260.


Satoshi Nakamoto, Bitcoin: A peer-to-peer electronic cash system (2008), https://bitcoin.org/ bitcoin.pdf


Leslie Lamport, Robert Shostak and Marshall Peaset, The Byzantine Generals Problem, ACM Transactions on Programming Languages and Systems Vol. 4, No. 3 (1982), pp. 382-401


Vitalik Buterin. Ethereum: A Next-Generation Smart Contract and Decentralized Application Platform (2013), https://github.com/ethereum/wiki/wiki/White-Paper


Delegated Proof-of-Stake Consensus, https://bitshares.org/technology/delegated-proof-of-stake-consensus/


Alan Demers, Dan Greene, Carl Hauser, Wes Irish, John Larson, Scott Shenker, Howard Sturgis, Dan Swinehart and Doug Terry, Epidemic algorithms for replicated database maintenance, PODC '87 Proceedings of the sixth annual ACM Symposium on Principles of distributed computing (1987), pp. 1-12.


Hyperledger Architecture, Vol. 1 (2017), https://www.hyperledger.org/wp-content/uploads/2017/08/Hyperledger Arch WG Paper 1 Consensus.pdf


Leemon Baird, The Swirlds Hashgraph Consensus Algorithm: Fair, Fast, Byzantine Fault Tolerance (2016), https://www.swirlds.com/downloads/SWIRLDS-TR-2016-01.pdf


David Mazirez, The Stellar Consensus Protocol: A Federated Model for Internet- level Consensus, Stellar Development Foundation (2016), https://www.stellar.org/papers/ stellar-consensus-protocol.pdf.


Gabriel Bracha and Sam Toueg, Asynchronous Consensus and Broadcast Protocols, Journal of the ACM, Vol. 32, Issue 4 (1985), pp. 824-840.


Miguel Castro and Barbara Liskov, Practical byzantine fault tolerance, Proceedings of the 3rd Symposium on Operating Systems Design and Implementation (1999), pp.173186.


Leslie Lamport, The Part-Time Parliament ACM Transactions on Computer Systems, Vol. 16, Issue 2 (1999), pp. 133-169.


Don Johnson, Alfred Menezes and Scott Van-Stone, The Elliptic Curve Digital Signature Algorithm (ECDSA), International Journal of Information Security, Vol. 1, Issue 1 (2001), pp. 36-63.


David L. Mills, Internet Time Synchronization: The Network Time Protocol, I.E.E.E. Trans. Comm. 39, No 10 (1991), pp. 1482-1493.


\subsection{https://blog.z.cash/snark-explain/}

\subsection{https://bitsofpy.blogspot.com/2014/03/homomorphic-encryption-using-rsa.html}

\subsection{https://bravenewcoin.com/assets/Whitepapers/QRL-whitepaper.pdf}

\end{document}
